% ============================================================
%  《材料制备与性能分析综合实践》课程报告
%  编译: xelatex 一阶段-daily-log-template.tex
% ============================================================

\documentclass[a4paper, 12pt]{ctexart}

% --- 宏包设置 ---
\usepackage{geometry}
\geometry{left=2.5cm, right=2.5cm, top=3cm, bottom=3cm}
\usepackage{tabularx}
\usepackage{ulem}
\usepackage{setspace}
\usepackage{graphicx}
\usepackage{subcaption}
\usepackage{float}
\usepackage{amsmath, amssymb}
\usepackage{booktabs}
\usepackage{xcolor}
\usepackage{hyperref}

% --- 页眉页脚 ---
\usepackage{fancyhdr}
\pagestyle{fancy}
\fancyhf{}
\fancyhead[L]{《材料制备与性能分析综合实践》课程报告}
\fancyhead[R]{\today}
\fancyfoot[C]{\thepage}
\renewcommand{\headrulewidth}{0.4pt}
\setlength{\headheight}{14.49998pt}
\addtolength{\topmargin}{-2.49998pt}

\begin{document}

% ==================== 封面页 ====================
\begin{titlepage}
    \vspace{6cm}

    % 主标题
    \begin{center}
        {\zihao{1} \heiti 《材料制备与性能分析综合实践》} \\[1cm]
        {\zihao{1} \heiti 课程报告}
    \end{center}

    \vspace{4cm}

    % 个人信息表单
    \begin{center}
        \zihao{3} \songti
        \renewcommand{\arraystretch}{1.8}
        \begin{tabular}{lc}
            \makebox[4em][c]{题目} & \underline{\makebox[8cm]{SIZrOC纤维表面改性探索}} \\
            \makebox[4em][c]{学院} & \underline{\makebox[8cm]{物理与材料科学学院}} \\
            \makebox[4em][c]{专业} & \underline{\makebox[8cm]{材料科学与工程}} \\
            \makebox[4em][c]{班级} & \underline{\makebox[8cm]{材科232(创)}} \\
            \makebox[4em][c]{姓名} & \underline{\makebox[8cm]{欧海轩}} \\
            \makebox[4em][c]{学号} & \underline{\makebox[8cm]{32319600008}} \\
        \end{tabular}
    \end{center}
\end{titlepage}

% ==================== 正文部分 ====================
\newpage
\setcounter{page}{1}
\section*{课程报告正文}

\subsection*{1. 引言}
SiZrOC 陶瓷纤维因其优异的高温稳定性与抗氧化性能，在航空航天及耐高温复合材料领域具有广泛的应用前景。本实践项目主要围绕 SiZrOC 纤维的表面改性展开探索。针对单一成分纤维在结构与力学性能上的局限性，本研究尝试引入多层静电纺丝技术，旨在通过控制不同层级的物理化学性质，优化纤维的前驱体成形质量及烧结后的界面结合特性。

\subsection*{2. 实验设计与制备}
在多层纺丝工艺的探索中，本实验拟定了两种实现方案：微观层面的同轴套管样式与宏观层面的"夹心饼干"叠层样式。受限于同轴挤出针头的硬件条件，本阶段主要采用宏观三层叠层方案进行实验。

实验采用有机层与无机层交替的方法进行表层成形。为维持有机网络的均匀性，在有机前驱体中加入了 1.72\,g 至 1.735\,g 的 PVP 作为纺丝助剂。在优化后的三层制备方案中，首底两层均采用含 1.735\,g PVP 的材料，而中间层则额外添加了 2\,g 葡萄糖与 0.64\,g 正丙醇锆。这种组分调节使得中间层表现出较强的粘性与一定的韧性。前驱体成形后，按照预设的程序（200/400/420/1250/60/−121）进行高温烧结处理。

\subsection*{3. 结果与讨论}
在多层前驱体的揭样剥离过程中，通过镊子可以轻易地将三层结构分离，且未观察到明显的粘滞拖拽现象，这表明各层之间保持了良好的独立性，层间界面清晰。烧结后产物的宏观形貌与前期单层样品无显著差异。

为了进一步观察纤维的微观形貌与表面改性效果，对样品进行了 SEM（扫描电子显微镜）检测。然而初步测试结果并不理想，经分析认为，测试环境（107 实验室）湿度偏高是导致形貌缺陷或成纤不良的主要原因。随后，实验设备被整体搬迁至环境控制更佳的创新大楼 A 座 201 房间，并重新完善了静电纺丝机及高压电源的接地等电气保障措施，为后续的 5321 硅烷原液改性实验（配比：15+15+0.96）奠定了基础。

\subsection*{4. 结论}
本研究初步验证了基于"夹心饼干"模式的三层静电纺丝工艺在 SiZrOC 纤维制备中的可行性。通过添加 PVP、葡萄糖与正丙醇锆，成功调节了中间层的粘韧性，实现了界面清晰的多层前驱体制备。后续实验将在优化后的低湿度环境中，进一步探究硅烷偶联剂对纤维界面的深度改性作用。

\newpage

% ==================== 日记部分 ====================
\begin{center}
    {\zihao{3} \heiti 《材料制备与性能分析综合实践》课程日记}
\end{center}
\vspace{0.5cm}

% 日记抬头信息
\noindent
\textbf{姓名：}欧海轩 \hfill \textbf{学号：}32319600008 \hfill \textbf{班级：}材科232(创)
\vspace{1cm}

% ============================================================
%  日记1：6月8日（周一）
% ============================================================
\subsection*{日记1：2026年6月8日 \quad 星期一}

\begin{itemize}
  \item 开始尝试新思路：多层纺丝。拟定两种实现方案：
  \begin{enumerate}
    \item 宏观层面——"夹心饼干"样式叠层
    \item 微观层面——同轴套管样式（需采购同轴挤出针头）
  \end{enumerate}
  \item 同轴针头价格较高，待与杜老师商议是否批准采购
  \item 配套工作：供液系统需改造，以解决单一电压点与双供液兼容问题，预计耗时较长
\end{itemize}

\vspace{1.5cm}

% ============================================================
%  日记2：6月9日（周二）
% ============================================================
\subsection*{日记2：2026年6月9日 \quad 星期二}

\begin{itemize}
  \item 6.2 无机的表层成形——纺丝
  \item 6.2 有机中加入 1.72\,g PVP，维持有机网络均匀性
  \item 中间层纺织（6.2 有机）
  \item 表层成型（6.2 无机）
\end{itemize}

\vspace{1.5cm}

% ============================================================
%  日记3：6月10日（周三）
% ============================================================
\subsection*{日记3：2026年6月10日 \quad 星期三}

\begin{itemize}
  \item 揭样（6.2 无机 / 6.2 有机 / 6.2 无机）—— 6.10 前驱体
  \item 烧结参数：200/400/420/1250/60/−121（6月11日执行，6月12日记录）
\end{itemize}

\vspace{1.5cm}

% ============================================================
%  日记4：6月14日（周日）
% ============================================================
\subsection*{日记4：2026年6月14日 \quad 星期日}

\begin{itemize}
  \item 纺织（6.2 有机）
\end{itemize}

\vspace{1.5cm}

% ============================================================
%  日记5：6月15日（周一）
% ============================================================
\subsection*{日记5：2026年6月15日 \quad 星期一}

新的优化措施：采用三层制备方案。其中，1.735\,g PVP 的材料担任首底两层的材料，中间层由额外添加了 2\,g 葡萄糖与 0.64\,g 正丙醇锆的组织担任。

对单独的中间层进行纺织时，其物理性质表现为有一定的韧性和较强的粘性。在前驱体成形后，可以用镊子轻易地将三层结构剥离分开——剥离过程中没有表现出粘滞拖拽的现象，说明层间界面清晰、各层独立性良好。烧结后，产物与前样品没有明显区别。

\vspace{1.5cm}

% ============================================================
%  日记6：6月18日（周四）
% ============================================================
\subsection*{日记6：2026年6月18日 \quad 星期四}

\begin{itemize}
  \item 下午拿样品找杜老师做 SEM 检测，结果不理想
  \item 怀疑是 107 实验室湿度偏高影响测试结果
  \item 杜老师指挥我与另一位同学将仪器从 107 实验室搬迁至创新大楼 A 座 201 房间
\end{itemize}

\vspace{1.5cm}

% ============================================================
%  日记7：6月22日（周一）
% ============================================================
\subsection*{日记7：2026年6月22日 \quad 星期一}

\begin{itemize}
  \item 5321 硅烷原液配制（6.23a）
\end{itemize}

\vspace{1.5cm}

% ============================================================
%  日记8：6月24日（周三）
% ============================================================
\subsection*{日记8：2026年6月24日 \quad 星期三}

\begin{itemize}
  \item 配比：15+15+0.96（6.24a）
\end{itemize}

\vspace{1.5cm}

% ============================================================
%  日记9：6月25日（周四）
% ============================================================
\subsection*{日记9：2026年6月25日 \quad 星期四}

\begin{itemize}
  \item 参加实习动员会。对比其他公司：一是任务分配过于流水线化，缺乏自由度与创造空间；二是研发方向不够前沿。经查阅资料后，决定前往盛立新材料，学习传统金属材料如何切入电子行业，更重要的是看好盛立所在行业的前景，希望深入接触。
  \item 找邢未未老师办理创新大楼 A 座 201 房间的门禁授权
\end{itemize}

\vspace{1.5cm}

% ============================================================
%  日记10：6月27日（周六）
% ============================================================
\subsection*{日记10：2026年6月27日 \quad 星期六}

\begin{itemize}
  \item 安装静电纺丝机及高压电源的电气保障措施（接地等），设备调试完成
  \item 考虑到下周需进工厂实习，实验时间不足，设备暂时闲置待用
\end{itemize}

\vspace{1.5cm}

\end{document}
